\appendix

\section{Notes on Fitting ideals}\label{sec:Fitt_gen}

In this section, we briefly review the definition of Fitting ideals and the shift theory introduced in~\cite{Kata_05}.

Let $R$ be a commutative ring.
Note that we do not assume $R$ is noetherian.
This is because the coefficient rings arising from the arithmetic in this paper are not noetherian in general (e.g., $\hZ[[\hZ]]$).
In~\cite{Kata_05}, the author developed the shift theory only over noetherian rings.
In this section, we observe that the argument can be generalized to non-noetherian rings by imposing appropriate finiteness properties on modules.

\subsection{Fitting ideals}\label{ss:Fitt_defn}

We recall the definition of Fitting ideals.
See, e.g., Northcott~\cite[\S 3.1]{Nor76}.

\begin{defi}\label{defn:Fitt1}
Let $I$ be a finite set and $J$ a (not necessarily finite) set.
Let
\[
h: R^{\oplus J} \to R^{\oplus I}
\]
be an $R$-homomorphism, where $R^{\oplus I}$, $R^{\oplus J}$ denote the free modules on the set $I$, $J$ respectively.

We define the (initial) Fitting ideal $\Fitt(h) = \Fitt_R(h)$ as follows.
For each subset $J' \subset J$ with $\# J' = \# I$, we write 
\[
h_{J'}: R^{\oplus J'} \to R^{\oplus I}
\]
for the restriction of $h$.
Let $\det(h_{J'})$ be the determinant of $h_{J'}$ with respect to any choice of bases; inevitably this determinant has ambiguity up to $R^{\times}$, but this does not matter.
Then $\Fitt_R(h)$ is defined as the ideal of $R$ generated by $\det(h_{J'})$ for all $J' \subset J$ with $\# J' = \# I$.
Note that if $\# J < \# I$, we have $\Fitt_R(h) = 0$ as there is not such a~$J'$.

More generally, for an integer $i \geq 0$, the $i$-th Fitting ideal $\Fitt_i(h) = \Fitt_{i, R}(h)$ is defined as follows.
For subsets $I' \subset I$ and $J' \subset J$ with $\# I' = \# J'$, we write 
\[
h_{J', I'}: R^{\oplus J'} \to R^{\oplus I'}
\]
for the map induced by $h$.
If $i \leq \# I$, we define $\Fitt_{i, R}(h)$ as the ideal of $R$ generated by $\det(h_{J', I'})$ for all $J' \subset J$, $I' \subset I$ with $\# J' = \# I' = \# I - i$.
If $i > \# I$, we set $\Fitt_{i, R}(h) = R$.
\end{defi}

\begin{defi}\label{defn:Fitt2}
Let $M$ be a finitely generated $R$-module.
We define $\Fitt_R(M)$ and $\Fitt_{i, R}(M)$ ($i \geq 0$) respectively as $\Fitt_R(h)$ and $\Fitt_{i, R}(h)$, where $h$ is a homomorphism as in Definition ~\ref{defn:Fitt1} such that the cokernel of $h$ is isomorphic to $M$.
It is known that this definition is independent from the choice of $h$.
\end{defi}

\subsection{Shifts of Fitting ideals}\label{ss:shift_defn}

We write $\Frac(R)$ for the total ring of fractions of~$R$.
An $R$-module $M$ is said to be torsion if any element is annihilated by a non-zero-divisor of $R$; equivalently if $\Frac(R) \otimes_R M = 0$.

\begin{lemm}\label{lem:tors_gen}
Let $F$ be a free $R$-module of finite rank and $h: F \to F$ be an endomorphism.
Then the following are equivalent.
\begin{itemize}
\item[(i)]
The map $h$ is injective.
\item[(ii)]
The determinant $\det(h) \in R$ is a non-zero-divisor.
\item[(iii)]
The cokernel $\Coker(h)$ is a torsion $R$-module.
\end{itemize}
\end{lemm}

\begin{proof}
The fact (i) $\Leftrightarrow$ (ii) is shown in~\cite[Chapter III, \S 8, Proposition 3, p.~524]{BouAlg74}.
Since $\Cok(h)$ is annihilated by $\det(h)$, we have (ii) $\Rightarrow$ (iii).
If (iii) holds, the base change of $h$ from $R$ to $\Frac(R)$ is surjective, so it is isomorphic, which implies (i).
\end{proof}

Let us define $\cP_R$ as the category of finitely presented torsion $R$-module $P$ such that $\pd_R(P) \leq 1$, where $\pd_R(-)$ denotes the projective dimension.
By Lemma ~\ref{lem:tors_gen}, if a module $P$ satisfies an exact sequence of the form
\begin{equation}\label{eq:pres}
0 \to F \to F \to P \to 0
\end{equation}
with $F$ a free module of finite rank, then $P$ is in $\cP_R$.
Conversely, if $R$ is local, then every $P$ in $\cP_R$ satisfies an exact sequence of the form \eqref{eq:pres}, since in that case any projective module is necessarily free.

A fractional ideal $I$ of $R$ is defined as an $R$-submodule of $\Frac(R)$ such that $u I \subset R$ for some non-zero-divisor $u \in R$.

\begin{lemm}\label{lem:FittP}
The following are true.
\begin{itemize}
\item[(1)]
For each $P \in \cP_R$, the Fitting ideal $\Fitt_R(P)$ is invertible as a fractional ideal of $R$.
\item[(2)]
Let $0 \to M' \to M \to P \to 0$ be an exact sequence of finitely generated torsion $R$-modules such that $P \in \cP_R$.
Then we have
\[
\Fitt_R(M) = \Fitt_R(P) \Fitt_R(M').
\]
\end{itemize}
\end{lemm}

\begin{proof}
We sketch the proof (see~\cite[Proposition 2.7]{Kata_05} for the details).
For both claims (1) and (2), it is enough to show them after localization at all prime ideals of $R$, so we may assume $R$ is local.
Then $P$ satisfies an exact sequence of the form \eqref{eq:pres}.
Now claim (1) follows from Lemma ~\ref{lem:tors_gen} and claim (2) follows by using the horseshoe lemma.
\end{proof}

Now we introduce the ``$n$-th shift'' $\Fitt^{[n]}_R(-)$ as in~\cite[Theorem 2.6]{Kata_05}.
First we consider the ``once-shift'' $\Fitt^{[1]}_R(-)$, which actually suffices for the purpose of this paper.
For completeness, in Definition ~\ref{defn:shift_n}, we will also introduce the general $n$-th shift.

\begin{defi}
For a finitely presented torsion $R$-module $M$, we define a fractional ideal $\Fitt^{[1]}_R(M)$ as follows.
Let us take an exact sequence of $R$-modules
\begin{equation}\label{eq:resol1}
0 \to N \to P \to M \to 0
\end{equation}
with $P \in \cP_R$.
For instance, if $M$ is generated by $n$ elements and annihilated by a non-zero-divisor $f \in R$, we may take $P = (R/fR)^n$.
Since $M$ is finitely presented, $N$ is finitely generated, so its Fitting ideal is defined.
Then we define (using Lemma ~\ref{lem:FittP}(1))
\[
\Fitt^{[1]}_R(M) = \Fitt_R(P)^{-1} \Fitt_R(N).
\]
This is well-defined, that is, independent from the choice of \eqref{eq:resol1}.
\end{defi}

Let us sketch the proof of the independency from \eqref{eq:resol1}.
Let $0 \to N' \to P' \to M \to 0$ be another exact sequence with $P' \in \cP_R$.
Then, by using the pull-back $L$ of the maps $P \to M$ and $P' \to M$, we obtain a commutative diagram with exact rows and columns
\[
\xymatrix{
	& & N' \ar@{=}[r] \ar@{^(->}[d]
	& N' \ar@{^(->}[d]
	& \\
	0 \ar[r]
	& N \ar[r] \ar@{=}[d]
	& L \ar[r] \ar@{->>}[d]
	& P' \ar[r] \ar@{->>}[d]
	& 0\\
	0 \ar[r]
	& N \ar[r]
	& P \ar[r]
	& M \ar[r]
	& 0.
}
\]
Since $P, P' \in \cP_R$, by Lemma ~\ref{lem:FittP}(2), we obtain
\[
\Fitt_R(L) = \Fitt_R(P') \Fitt_R(N),
\quad
\Fitt_R(L) = \Fitt_R(P) \Fitt_R(N').
\]
These formulas imply 
\[
\Fitt_R(P)^{-1} \Fitt_R(N)
= \Fitt_R(P')^{-1} \Fitt_R(N'),
\]
as desired.

Finally let us introduce the general $n$-th shift.
See~\cite[Theorem 2.6]{Kata_05} for the proof; to remove the noetherian hypothesis, we only have to suitably deal with the finiteness conditions.

\begin{defi}\label{defn:shift_n}
Let $n \geq 0$ be an integer.
Let $M$ be a torsion $R$-module such that there exists an exact sequence
\[
R^{a_n} \to R^{a_{n-1}} \to \dots \to R^{a_0} \to M \to 0
\]
for some integers $a_0, \dots, a_n \geq 0$.
For instance, when $n = 0$ (resp.~$n = 1$), this condition means that $M$ is finitely generated (resp.~finitely presented).
For such an $M$, we define a fractional ideal $\Fitt^{[n]}_R(M)$ as follows.
Let us take an exact sequence of $R$-modules
\begin{equation}\label{eq:resol2}
0 \to N \to P_1 \to \cdots \to P_n \to M \to 0
\end{equation}
with $P_1, \dots, P_n \in \cP_R$.
The existence of such a sequence follows from the assumption on $M$, and moreover then $N$ is finitely generated.
Then we define (using Lemma ~\ref{lem:FittP}(1))
\[
\Fitt^{[n]}_R(M) = \left( \prod_{i=1}^n \Fitt_R(P_i)^{(-1)^i} \right) \Fitt_R(N).
\]
This is well-defined, that is, independent from the choice of \eqref{eq:resol2}.
\end{defi}
 
